\chapter{Совместный гессиан singlet-вакуума и виртуального цикла}
\label{ch:v7-singlet-vacuum-virtual-cycle-combined-hessian-gate}

\section{Минимальная безразмерная модель}

Gauge-Casimir селектор запускает четыре бесцветных ребра. Два из них,
$p=z_{X_Le_R}$ и $q=z_{L_LY_R}$, входят в виртуальный циклический
детерминант; две вектороподобные массы не входят и сохраняют исходный
Hodge-потенциал.

После измерения амплитуд $p,q$ в единицах Hodge-уровня $\mu$ положим
\begin{equation}
 r=\frac{|p|}{\mu},\qquad s=\frac{|q|}{\mu},\qquad
 a=\left(\frac{\kappa\mu^2}{M_aM_b}\right)^2.
 \label{eq:v7-combined-hessian-variables}
\end{equation}
Конечномерный контроль совместного потенциала имеет вид
\begin{equation}
 V(r,s)=(r^2-1)^2+(s^2-1)^2
 +\gamma\log(1-ar^2s^2),
 \qquad ar^2s^2<1,
 \label{eq:v7-combined-hessian-potential}
\end{equation}
где $\gamma>0$ есть относительный вес determinant-блока. Параметр $a$
измеряет смешивание в единицах тяжёлой щели. На этом этапе оба числа
сохраняются явными: предыдущие гейты не вывели их отношение одним
спектральным профилем.

\section{Симметричная стационарная ветвь}

Искомая ветвь имеет
\begin{equation}
 r=s=\sqrt u,\qquad u>1.
 \label{eq:v7-combined-hessian-symmetric-branch}
\end{equation}
Стационарность фиксирует не $u$ отдельно, а связь параметров
\begin{equation}
 \gamma=\frac{2(u-1)(1-au^2)}{au}.
 \label{eq:v7-combined-hessian-stationarity}
\end{equation}
Положительность $\gamma$ требует открытой тяжёлой щели
\begin{equation}
 au^2<1.
 \label{eq:v7-combined-hessian-gap}
\end{equation}
Детерминант сдвигает радиус наружу, $u>1$, но не заставляет цветные поля
получать ненулевое среднее.

\section{Точные радиальные собственные значения}

В симметричной точке два связанных радиальных направления распадаются на
синфазное и противофазное. После использования условия стационарности их
собственные значения равны
\begin{align}
 \lambda_{\parallel}
 &=8\frac{1-au^2(2u-1)}{1-au^2},
 \label{eq:v7-combined-hessian-radial-parallel}\\
 \lambda_{\perp}&=8(2u-1).
 \label{eq:v7-combined-hessian-radial-perpendicular}
\end{align}
Второе значение положительно для всей рассматриваемой ветви. Первое
положительно тогда и только тогда, когда
\begin{equation}
 \boxed{au^2(2u-1)<1.}
 \label{eq:v7-combined-hessian-local-stability}
\end{equation}
Это условие строго сильнее \eqref{eq:v7-combined-hessian-gap}; поэтому
локально устойчивый минимум автоматически лежит внутри области, где
цветная пара остаётся тяжёлой.

Например,
\begin{equation}
 a=\frac1{10},\qquad u=\frac65,qquad
 \gamma=\frac{214}{75}
 \label{eq:v7-combined-hessian-benchmark}
\end{equation}
даёт положительные $\lambda_{\parallel}$ и $\lambda_{\perp}$ и ненулевую
минимальную тяжёлую массу. Это сертификат непустоты устойчивой области, а
не физическая калибровка параметров.

\section{Полная структура гессиана}

Две вектороподобные массовые амплитуды не входят в determinant-слово и
сохраняют радиальные собственные значения $8$ в выбранных безразмерных
единицах. Потенциал зависит только от модулей $p,q$, поэтому четыре фазы
четырёх singlet-рёбер остаются нулевыми модами:
\begin{equation}
 (n_-,n_0,n_+)_{\rm singlet}=(0,4,4)
 \label{eq:v7-combined-hessian-singlet-signature}
\end{equation}
в области \eqref{eq:v7-combined-hessian-local-stability}. Все пять
запрещённых и два виртуальных цветных блока сохраняют положительные
квадратичные массы; последние дополнительно требуют
\eqref{eq:v7-combined-hessian-gap}. Следовательно, полный локальный
гессиан не получает отрицательных мод.

Однако $|pq|^2$ не зависит от относительной фазы $p/q$. Виртуальный
детерминант не фиксирует фазовую или семейную ориентацию и не переоткрывает
замороженный CKM/PMNS-маршрут.

\section{Глобальная и нормировочная граница}

В конечномерной модели
\begin{equation}
 \lim_{ar^2s^2\to1^-}V(r,s)=-\infty.
 \label{eq:v7-combined-hessian-boundary-divergence}
\end{equation}
Поэтому найденная точка является только локальным метастабильным минимумом:
нулевая тяжёлая мода появляется на границе, где гауссово исключение теряет
применимость. В полном четырёхмерном эффективном потенциале поведение должно
быть вычислено после перенормировки; перенос конечномерного логарифма на
глобальный Wilson-потенциал запрещён. Локальные и глобальные критерии
устойчивости многополевого эффективного потенциала необходимо различать
\cite{v7-combined-hessian-cw,v7-combined-hessian-kannike}.

Кроме того, стационарность \eqref{eq:v7-combined-hessian-stationarity}
связывает два пока свободных безразмерных числа $a$ и $\gamma$. Наличие
открытой устойчивой области не выбирает одну физическую точку.

\begin{theorem}[Условный локальный проход совместного вакуума]
Существует непустая открытая область параметров, в которой четыре
gauge-синглетных ребра имеют стационарный Hodge--determinant вакуум,
совместный радиальный гессиан положителен, а два цветных циклических моста
остаются тяжёлыми с нулевыми средними. Тем самым локальная совместимость
singlet-запуска и виртуального цикла доказана. Однако конечномерный
потенциал не имеет доказанного глобального минимума, а параметры $a$ и
$\gamma$ не выведены одним действием.
\end{theorem}

\begin{nogocriterion}[Запрет повышения локального минимума до вакуума]
Нельзя объявлять найденную устойчивую точку физическим вакуумом только по
локальному гессиану. Требуются четырёхмерный перенормированный потенциал,
глобальная ограниченность или контролируемое время жизни, а также вывод
$a$ и $\gamma$ из одного спектрального профиля до выбора benchmark.
\end{nogocriterion}

\begin{protocol}
\begin{enumerate}
 \item цветосохраняющий singlet-Hodge вакуум: \textbf{использован};
 \item тяжёлая область: \textbf{$ar^2s^2<1$};
 \item симметричная стационарная ветвь: \textbf{получена};
 \item условие локальной устойчивости: \textbf{$au^2(2u-1)<1$};
 \item устойчивая область непуста: \textbf{да};
 \item отрицательные локальные моды: \textbf{$0$};
 \item фазовые нули singlet-сектора: \textbf{$4$};
 \item цветная тяжёлая щель в устойчивой области: \textbf{открыта};
 \item фиксация относительной фазы $p/q$: \textbf{нет};
 \item глобальный конечномерный минимум: \textbf{нет};
 \item четырёхмерный перенормированный потенциал: \textbf{не вычислен};
 \item параметры $a,\gamma$ выведены одним действием: \textbf{нет};
 \item итог: \textbf{условный локальный проход};
 \item следующий гейт: \textbf{общий спектральный профиль отношения
 singlet--virtual}.
\end{enumerate}
\end{protocol}

Машинный сертификат сохранён в
\path{s2t_v7_singlet_vacuum_virtual_cycle_combined_hessian_gate_results.json}.

\begin{thebibliography}{9}
\bibitem{v7-combined-hessian-cw}
S.~Coleman, E.~Weinberg,
\emph{Radiative Corrections as the Origin of Spontaneous Symmetry Breaking},
Physical Review D 7 (1973), 1888--1910.
\bibitem{v7-combined-hessian-kannike}
K.~Kannike,
\emph{Vacuum Stability of a General Scalar Potential of a Few Fields},
European Physical Journal C 76 (2016), 324; arXiv:1603.02680.
\end{thebibliography}